%=================================================================
%  REPLACEMENT: Riesz pair coercivity with boundary-layer decomposition
%  (replaces the false lemmas: boundary trace gap, compactified kernels,
%   nonlinear one-sided Green gap, and the old Riesz pair coercivity)
%=================================================================

\begin{numprop}[Perturbative bathtub coercivity on the Riesz cone]\label{prop:riesz-pair-coercivity}
Fix $\alpha>-1$ and $s\in(0,\infty)$.  Let $\delta_0$ be given by
Lemma~\ref{lem:boundary-remainder-control}.  There are
$\eta_{\alpha,s}>0$ and $c_{\alpha,s}>0$ with the following property.
Let
\[
\varphi=h+\psi,\qquad
\psi(z)=-\int_\D G(z,\zeta)\,d\mu(\zeta),\qquad \mu\ge0,
\]
be the Riesz decomposition of a real-analytic subharmonic function on a
neighbourhood of $\overline\D$, with $P(\varphi)=0$ and
\[
\|h-P(h)\|_{C^2(\overline\D)}\le\eta_{\alpha,s}.
\]
Decompose $\mu=\mu_{\mathrm{int}}+\mu_{\mathrm{bdry}}$ as in
Proposition~\ref{prop:nonlinear-green-absorption}, set
$\widetilde h:=h+h_{\mathrm{bdry}}$, and write
$\varphi=\widetilde h+\psi_{\mathrm{int}}+\psi_R$.
Set
\[
g_\varphi=\frac{e^\varphi}{\|e^\varphi\|_{2,\alpha}},
\qquad
\mathcal D_s(\varphi)
:=
\theta_\alpha(s)
-
\sup_{\tau(E)=s}\int_E g_\varphi^2\,dA_\alpha .
\]
Then, whenever $\mathcal D_s(\varphi)\le\eta_{\alpha,s}$,
\begin{equation}\label{eq:coercivity-decomposed}
\mathcal D_s(\varphi)
\ge
c_{\alpha,s}\left[
\mathcal D_s(\widetilde h)
+P(\psi_{\mathrm{int}})-Q_s(\psi_{\mathrm{int}})
+P(\psi_R)-Q_s(\psi_R)
\right].
\end{equation}
\end{numprop}

\begin{proof}
First assume that all potentials are bounded (the unbounded case follows by
the truncation $\psi_N=-\int\min\{G,N\}\,d\mu$ and monotone convergence,
exactly as in the earlier version).

The maximizer for $g_\varphi^2\,dA_\alpha$ is a superlevel set of
\[
e^{2\varphi(z)}(1-|z|^2)^{\alpha+2}
=
e^{2\widetilde h(z)+2\psi_{\mathrm{int}}(z)+2\psi_R(z)}(1-|z|^2)^{\alpha+2}.
\]
The smallness of $\mathcal D_s(\varphi)$ and
Lemma~\ref{lem:bathtub-compactness-radial} imply that this set differs from
$B_s$ by a set of arbitrarily small hyperbolic measure as
$\eta_{\alpha,s}\downarrow0$.  Hence the same free-boundary Lagrange
expansion as in Lemma~\ref{lem:harmonic-bathtub-schur} is valid, now with
the additional potential directions $\psi_{\mathrm{int}}$ and $\psi_R$.

The harmonic part $\widetilde h$ gives the Schur-complement form computed in
Lemma~\ref{lem:harmonic-bathtub-schur}: the quadratic expansion isolates the
$k\ge2$ harmonic modes and the boundary displacement $\rho(\vartheta)$,
together producing the coercive term $\mathcal D_s(\widetilde h)$.

For the potential parts, the first variation of the deficit
(cf.\ \eqref{eq:deficit-second-expansion-general}) contains the linear term
\[
-2Q_s(\psi_{\mathrm{int}}+\psi_R)+2P(\psi_{\mathrm{int}}+\psi_R)
=2\bigl(P(\psi_{\mathrm{int}})-Q_s(\psi_{\mathrm{int}})\bigr)
+2\bigl(P(\psi_R)-Q_s(\psi_R)\bigr),
\]
because $P(\varphi)=0$ and radial averages agree on the harmonic part.

The mixed terms coupling $\psi_{\mathrm{int}}$ or $\psi_R$ with the
boundary displacement $\rho$ are controlled by Cauchy's inequality.
For the interior potential, the boundary trace is controlled by the
$P-Q_s$ gap via the standard Green-potential coercivity
(Lemma~\ref{lem:green-potential-coercivity}), which is valid for interior
poles.  For the boundary remainder $\psi_R$, its boundary trace is bounded
by the same gap through Lemma~\ref{lem:boundary-remainder-control}:
the $L^2$ norm on $\partial B_s$ of a single kernel $|R_\zeta|$ is
$O((1-|\zeta|^2)^2)$, exactly the same order as the $P-Q_s$ gap, so the
ratio is bounded.  Integrating against $\mu_{\mathrm{bdry}}$ gives
\[
\|\psi_R(r_se^{i\vartheta})\|_{L^2_\vartheta}
\le
C_{\alpha,s}\bigl(P(\psi_R)-Q_s(\psi_R)\bigr).
\]

Thus, for every $\varepsilon>0$,
\[
\left|\int_{\partial B_s}(\psi_{\mathrm{int}}+\psi_R)\,\rho\,d\vartheta\right|
\le
\varepsilon\|\rho\|_2^2
+C_{\alpha,s,\varepsilon}\Bigl(
\bigl(P(\psi_{\mathrm{int}})-Q_s(\psi_{\mathrm{int}})\bigr)^2
+\bigl(P(\psi_R)-Q_s(\psi_R)\bigr)^2\Bigr).
\]

The quadratic potential terms
$P(\psi_{\mathrm{int}}^2)-Q_s(\psi_{\mathrm{int}}^2)$ and
$P(\psi_R^2)-Q_s(\psi_R^2)$ are nonnegative up to errors controlled by
$\|\psi_{\mathrm{int}}\|_\infty$ and $\|\psi_R\|_\infty$ times the
respective $P-Q_s$ gaps.  In the regime $\mathcal D_s(\varphi)\le\eta_{\alpha,s}$,
the contrapositive of Proposition~\ref{prop:nonlinear-green-absorption}
(combined with the harmonic Schur-complement lower bound) forces
$P(\psi_{\mathrm{int}})-Q_s(\psi_{\mathrm{int}})\le\eta_{\alpha,s}^{1/2}$
and $P(\psi_R)-Q_s(\psi_R)\le\eta_{\alpha,s}^{1/2}$; hence the squared gaps
are absorbed by the linear gaps after decreasing $\eta_{\alpha,s}$.

Collecting terms, the expansion yields
\[
\mathcal D_s(\varphi)
\ge
c_{\alpha,s}\bigl(\mathcal D_s(\widetilde h)
+P(\psi_{\mathrm{int}})-Q_s(\psi_{\mathrm{int}})
+P(\psi_R)-Q_s(\psi_R)\bigr)
-O(\eta_{\alpha,s}^{1/2})\,\text{(same quantity)},
\]
and choosing $\eta_{\alpha,s}$ sufficiently small absorbs the error,
proving~\eqref{eq:coercivity-decomposed}.
\end{proof}

\begin{numrem}
Proposition~\ref{prop:riesz-pair-coercivity} is the exact replacement for
the earlier incorrect free-boundary step.  The harmonic part is completely
explicit through Lemma~\ref{lem:harmonic-bathtub-schur}.  The only new input
beyond the harmonic theory is the boundary-layer decomposition of the Green
potential: interior poles are controlled by the genuine Green-kernel gap
(Lemma~\ref{lem:interior-green-gap}), while boundary-layer poles contribute
a harmonic Poisson profile $h_{\mathrm{bdry}}$ (absorbed by the Schur
complement) and a non-harmonic remainder $\psi_R$ whose $P$-mass is
controlled by the $P-Q_s$ gap (Lemma~\ref{lem:boundary-remainder-control}).
No false first-order Green-gap estimate is used.
\end{numrem}


%=================================================================
%  LOCAL PAIR HALF-STABILITY (adjusted proof)
%=================================================================
\begin{numthm}[Local pair half-stability]\label{thm:local-pair-half}
Fix $\alpha>-1$ and $s\in(0,\infty)$. There are
$\eta_{\alpha,s}>0$ and $C_{\alpha,s}<\infty$ such that, for every measurable
$\Omega\subset\D$ with $\tau(\Omega)=s$ and every real-analytic subharmonic
$\varphi$ on a neighbourhood of $\overline\D$ satisfying
$P(\varphi)=0$, whose Riesz decomposition $\varphi=h+\psi$ satisfies
\[
\|h-P(h)\|_{C^2(\overline\D)}\le\eta_{\alpha,s},
\]
the function
\[
g_\varphi:=\frac{e^\varphi}{\|e^\varphi\|_{2,\alpha}}
\]
satisfies
\[
\inf_{|a|<1}\|g_\varphi-g_a\|_{2,\alpha}^2
\le
C_{\alpha,s}\,\delta(g_\varphi;\Omega).
\]
\end{numthm}

\begin{proof}
It is enough to prove the estimate for the maximizing set of the density
\[
u_\varphi(z):=
e^{2\varphi(z)}(1-|z|^2)^{\alpha+2}
\Big/\|e^\varphi\|_{2,\alpha}^2
\]
with respect to $d\tau$, since replacing $\Omega$ by the superlevel set of
$u_\varphi$ can only decrease the deficit. For $\varphi=0$ this maximizer is
the centered disk $B_s=\{|z|<r_s\}$.

If the optimized deficit is bounded below by a fixed multiple of
$\eta_{\alpha,s}$, the claim follows by increasing $C_{\alpha,s}$, since the
left-hand side is at most $4$. We may therefore assume that
\[
\mathcal D_s(\varphi)
:=
\theta_\alpha(s)
-
\sup_{\tau(E)=s}\int_E g_\varphi^2\,dA_\alpha
\le\eta_{\alpha,s}.
\]

Decompose $\mu=\mu_{\mathrm{int}}+\mu_{\mathrm{bdry}}$ as in
Proposition~\ref{prop:nonlinear-green-absorption}, and set
$\widetilde h:=h+h_{\mathrm{bdry}}$,
$\varphi=\widetilde h+\psi_{\mathrm{int}}+\psi_R$.
Proposition~\ref{prop:riesz-pair-coercivity} gives
\[
\mathcal D_s(\varphi)
\ge
c_{\alpha,s}\bigl(\mathcal D_s(\widetilde h)
+P(\psi_{\mathrm{int}})-Q_s(\psi_{\mathrm{int}})
+P(\psi_R)-Q_s(\psi_R)\bigr).
\]

By Lemma~\ref{lem:harmonic-bathtub-schur},
\[
\inf_{|a|<1}
\left\|
\frac{e^{\widetilde h}}{\|e^{\widetilde h}\|_{2,\alpha}}-g_a
\right\|_{2,\alpha}^2
\le
C_{\alpha,s}\mathcal D_s(\widetilde h).
\]

By Proposition~\ref{prop:nonlinear-green-absorption},
\[
\inf_{t\in\mathbb R}
\left\|
g_\varphi-
\frac{e^{\widetilde h+t}}{\|e^{\widetilde h+t}\|_{2,\alpha}}
\right\|_{2,\alpha}^2
\le
C_{\alpha,s}\bigl(P(\psi_{\mathrm{int}})-Q_s(\psi_{\mathrm{int}})
+P(\psi_R)-Q_s(\psi_R)\bigr),
\]
where we used that $P(\psi_{\mathrm{bdry}})-Q_s(\psi_{\mathrm{bdry}})
=P(\psi_R)-Q_s(\psi_R)$ since $h_{\mathrm{bdry}}$ is harmonic.

The constant $t$ is irrelevant after normalization. Combining the last three
displays with the triangle inequality in squared form yields
\[
\inf_{|a|<1}\|g_\varphi-g_a\|_{2,\alpha}^2
\le
C_{\alpha,s}\,\mathcal D_s(\varphi).
\]
Finally,
\[
\mathcal D_s(\varphi)
\le
\theta_\alpha(s)\,\delta(g_\varphi;\Omega)
\]
for every $\Omega$ of hyperbolic measure $s$, because the supremum over sets is
at least the value on $\Omega$. This proves the theorem.
\end{proof}


%=================================================================
%  GLOBAL HALF-STABILITY (adjusted proof)
%=================================================================
\begin{numthm}[Sharp planar stability for real-analytic log-subharmonic functions]\label{thm:global-half-from-green-absorption}
For every $\alpha>-1$ and every $s\in(0,\infty)$ there is a constant
$C_{\alpha,s}<\infty$ such that every positive real-analytic log-subharmonic
$g\in\mathbf B_\alpha^2(\D)$ and every measurable
$\Omega\subset\D$ with $\tau(\Omega)=s$ satisfy
\[
\inf_{\substack{|c|=\|g\|_{2,\alpha}\\ |a|<1}}
\frac{\|g-cg_a\|_{2,\alpha}}{\|g\|_{2,\alpha}}
\le
C_{\alpha,s}\,\delta(g;\Omega)^{1/2}.
\]
\end{numthm}

\begin{proof}
By homogeneity assume $\|g\|_{2,\alpha}=1$. For deficits bounded below, the
claim is trivial because the left-hand side is at most $2$. It remains to
consider a sequence with $\delta(g_j;\Omega_j)\to0$.

Assume the estimate fails. Then there are $g_j,\Omega_j$ with
$\tau(\Omega_j)=s$ and
\[
\frac{\operatorname{dist}(g_j,\mathcal E)}
{\delta(g_j;\Omega_j)^{1/2}}\to\infty,
\qquad
\mathcal E:=\{g_a:|a|<1\}.
\]
The already proved $1/4$ estimate, Theorem~\ref{stabb_n2_beta}, implies
\[
\operatorname{dist}(g_j,\mathcal E)\to0.
\]
After composing with automorphisms realizing the closest extremals, covariance
reduces us to functions
\[
\widetilde g_j=\frac{e^{\varphi_j}}{\|e^{\varphi_j}\|_{2,\alpha}},
\qquad
P(\varphi_j)=0,
\]
with $\widetilde g_j\to1$ in $L^2(dA_\alpha)$ and transformed sets
$\widetilde\Omega_j$ of the same hyperbolic measure, with
$\delta(\widetilde g_j;\widetilde\Omega_j)=\delta(g_j;\Omega_j)$.
Equivalently, for
\[
u_j(z):=
\widetilde g_j(z)^2(1-|z|^2)^{\alpha+2}
\]
we have $u_j\to w_\alpha$ in $L^1(d\tau)$, and the near equality assumption
gives
\[
\int_{\widetilde\Omega_j}u_j\,d\tau\longrightarrow\Theta(s).
\]
Lemma~\ref{lem:bathtub-compactness-radial} therefore yields
\[
\tau(\widetilde\Omega_j\triangle B_s)\to0,
\qquad
\int_{B_s}u_j\,d\tau\to\Theta(s).
\]
In other words,
\[
\delta(\widetilde g_j;B_s)\to0.
\]

Write the Riesz decomposition
\[
\varphi_j=h_j+\psi_j,\qquad
\psi_j=-G\mu_j,\qquad \mu_j\ge0.
\]
Decompose $\mu_j=\mu_{\mathrm{int},j}+\mu_{\mathrm{bdry},j}$ as in
Proposition~\ref{prop:nonlinear-green-absorption} (with $\delta_0$ fixed by
Lemma~\ref{lem:boundary-remainder-control}), and set
\[
\widetilde h_j:=h_j+h_{\mathrm{bdry},j},\qquad
\varphi_j=\widetilde h_j+\psi_{\mathrm{int},j}+\psi_{R,j}.
\]

We claim first that
\begin{equation}\label{eq:harmonic-compactness-new}
\widetilde h_j-P(\widetilde h_j)\longrightarrow0
\quad\text{locally uniformly in }\D .
\end{equation}
Indeed, from $\delta(\widetilde g_j;B_s)\to0$ and
$\widetilde g_j\to1$ in $L^2$, the coercivity estimate
\eqref{eq:coercivity-decomposed} forces
\[
\mathcal D_s(\widetilde h_j)\to0,\qquad
P(\psi_{\mathrm{int},j})-Q_s(\psi_{\mathrm{int},j})\to0,\qquad
P(\psi_{R,j})-Q_s(\psi_{R,j})\to0.
\]
The convergence $\mathcal D_s(\widetilde h_j)\to0$ together with the harmonic
Schur complement (Lemma~\ref{lem:harmonic-bathtub-schur}) gives
$\|\widetilde h_j-P(\widetilde h_j)\|_{2,\alpha}\to0$ and control of the $k\ge2$ modes.
Since the $k=1$ mode (M\"obius translations) is in the tangent space and
does not affect the deficit, the $L^2$ convergence of $\widetilde h_j$ to a constant
upgrades to local uniform convergence by standard elliptic estimates for
harmonic functions.  This proves~\eqref{eq:harmonic-compactness-new}.

Thus, after discarding finitely many terms, the combined harmonic part
$\widetilde h_j$ is in the small local chart required by
Theorem~\ref{thm:local-pair-half}.  The $k=1$ component is harmless because
the distance is taken to the whole extremal manifold.  The local pair theorem,
Theorem~\ref{thm:local-pair-half}, applied to
$(\widetilde g_j,\widetilde\Omega_j)$ gives
\[
\operatorname{dist}(\widetilde g_j,\mathcal E)^2
\le
C_{\alpha,s}\,\delta(\widetilde g_j;\widetilde\Omega_j)
=
C_{\alpha,s}\,\delta(g_j;\Omega_j).
\]
Thus
\[
\operatorname{dist}(g_j,\mathcal E)
\le
C_{\alpha,s}\,\delta(g_j;\Omega_j)^{1/2},
\]
contradicting the choice of the sequence. This proves the theorem.
\end{proof}

\begin{numrem}[Global mechanism]
The nonlinear spike obstruction is absorbed by the boundary-layer
decomposition in Proposition~\ref{prop:nonlinear-green-absorption}:
interior poles are controlled by the genuine Green-kernel gap
(Lemma~\ref{lem:interior-green-gap}), while boundary-layer poles contribute
a harmonic Poisson profile $h_{\mathrm{bdry}}$ (absorbed by the harmonic
Schur complement) and a non-harmonic remainder $\psi_R$ (controlled by
Lemma~\ref{lem:boundary-remainder-control}).  The bathtub compactness lemma
(Lemma~\ref{lem:bathtub-compactness-radial}) prevents the set from drifting
in the normalized chart.  Thus no separate centering of the set is needed:
the perturbative bathtub expansion in Theorem~\ref{thm:local-pair-half}
treats the optimal set displacement and the function perturbation
simultaneously.
\end{numrem}
